% Options for packages loaded elsewhere
\PassOptionsToPackage{unicode}{hyperref}
\PassOptionsToPackage{hyphens}{url}
\PassOptionsToPackage{dvipsnames,svgnames,x11names}{xcolor}
\documentclass[
  10pt,
]{article}
\usepackage{xcolor}
\usepackage[margin=1cm,top=1cm,bottom=1cm,left=1cm,right=1cm,includeheadfoot]{geometry}
\usepackage{amsmath,amssymb}
\setcounter{secnumdepth}{3}
\usepackage{iftex}
\ifPDFTeX
  \usepackage[T1]{fontenc}
  \usepackage[utf8]{inputenc}
  \usepackage{textcomp} % provide euro and other symbols
\else % if luatex or xetex
  \usepackage{unicode-math} % this also loads fontspec
  \defaultfontfeatures{Scale=MatchLowercase}
  \defaultfontfeatures[\rmfamily]{Ligatures=TeX,Scale=1}
\fi
\usepackage{lmodern}
\ifPDFTeX\else
  % xetex/luatex font selection
  \setmainfont[]{DejaVu Serif}
  \setmonofont[]{DejaVu Sans Mono}
\fi
% Use upquote if available, for straight quotes in verbatim environments
\IfFileExists{upquote.sty}{\usepackage{upquote}}{}
\IfFileExists{microtype.sty}{% use microtype if available
  \usepackage[]{microtype}
  \UseMicrotypeSet[protrusion]{basicmath} % disable protrusion for tt fonts
}{}
\usepackage{setspace}
\makeatletter
\@ifundefined{KOMAClassName}{% if non-KOMA class
  \IfFileExists{parskip.sty}{%
    \usepackage{parskip}
  }{% else
    \setlength{\parindent}{0pt}
    \setlength{\parskip}{6pt plus 2pt minus 1pt}}
}{% if KOMA class
  \KOMAoptions{parskip=half}}
\makeatother
\usepackage{listings}
\newcommand{\passthrough}[1]{#1}
\lstset{defaultdialect=[5.3]Lua}
\lstset{defaultdialect=[x86masm]Assembler}
\usepackage{graphicx}
\makeatletter
\newsavebox\pandoc@box
\newcommand*\pandocbounded[1]{% scales image to fit in text height/width
  \sbox\pandoc@box{#1}%
  \Gscale@div\@tempa{\textheight}{\dimexpr\ht\pandoc@box+\dp\pandoc@box\relax}%
  \Gscale@div\@tempb{\linewidth}{\wd\pandoc@box}%
  \ifdim\@tempb\p@<\@tempa\p@\let\@tempa\@tempb\fi% select the smaller of both
  \ifdim\@tempa\p@<\p@\scalebox{\@tempa}{\usebox\pandoc@box}%
  \else\usebox{\pandoc@box}%
  \fi%
}
% Set default figure placement to htbp
\def\fps@figure{htbp}
\makeatother
\setlength{\emergencystretch}{3em} % prevent overfull lines
\providecommand{\tightlist}{%
  \setlength{\itemsep}{0pt}\setlength{\parskip}{0pt}}
% Enable graphics inclusion and ensure figure numbering works
\usepackage{graphicx}
\renewcommand{\figurename}{Figure}

% Configure fonts for Unicode support with fallbacks
\usepackage{newunicodechar}
\newunicodechar{⁴}{\textsuperscript{4}}
\newunicodechar{₄}{\textsubscript{4}}

% Math notation macros
\newcommand{\norm}[1]{\lVert #1\rVert}
\newcommand{\abs}[1]{\lvert #1\rvert}

% Enhanced code block styling for better contrast and readability
\usepackage{fancyvrb}
\usepackage{xcolor}
\usepackage{listings}

% Define custom colors for code blocks
\definecolor{codebg}{RGB}{245, 245, 245}      % Light gray background
\definecolor{codeborder}{RGB}{200, 200, 200}  % Medium gray border
\definecolor{codefg}{RGB}{50, 50, 50}         % Dark gray text

% Configure Verbatim environment for inline code
\DefineVerbatimEnvironment{Verbatim}{Verbatim}{%
    fontsize=\small,
    frame=single,
    framerule=0.5pt,
    framesep=3pt,
    rulecolor=\color{codeborder},
    bgcolor=\color{codebg},
    fgcolor=\color{codefg}
}

% Configure code block styling
\DefineVerbatimEnvironment{Highlighting}{Verbatim}{%
    fontsize=\footnotesize,
    frame=single,
    framerule=0.5pt,
    framesep=5pt,
    rulecolor=\color{codeborder},
    bgcolor=\color{codebg},
    fgcolor=\color{codefg}
}

% Style inline code with \texttt
\renewcommand{\texttt}[1]{%
    \colorbox{codebg}{\color{codefg}\ttfamily #1}%
}

% Configure listings package for code blocks
\lstset{
    backgroundcolor=\color{codebg},
    basicstyle=\footnotesize\ttfamily\color{codefg},
    breakatwhitespace=false,
    breaklines=true,
    captionpos=b,
    commentstyle=\color{codefg},
    deletekeywords={...},
    escapeinside={\%*}{*)},
    extendedchars=true,
    frame=single,
    framerule=0.5pt,
    framesep=5pt,
    keepspaces=true,
    keywordstyle=\color{codefg},
    language=Python,
    morekeywords={*,...},
    numbers=left,
    numbersep=5pt,
    numberstyle=\tiny\color{codefg},
    rulecolor=\color{codeborder},
    showspaces=false,
    showstringspaces=false,
    showtabs=false,
    stepnumber=1,
    stringstyle=\color{codefg},
    tabsize=2,
    title=\lstname
}

% Override any Pandoc default lstset configurations
\AtBeginDocument{
    \lstset{
        backgroundcolor=\color{codebg},
        basicstyle=\footnotesize\ttfamily\color{codefg},
        frame=single,
        framerule=0.5pt,
        framesep=5pt,
        rulecolor=\color{codeborder},
        numbers=left,
        numbersep=5pt,
        numberstyle=\tiny\color{codefg}
    }
}

% Configure hyperref colors consistently
\AtBeginDocument{
% Override pandoc's hidelinks setting with consistent options
\hypersetup{
    colorlinks=true,
    allcolors=red,
    linkcolor=red,
    urlcolor=red,
    citecolor=red,
    filecolor=red,
    menucolor=red,
    linktoc=all
}
}

% Simple page break support for document structure
\usepackage{bookmark}
\IfFileExists{xurl.sty}{\usepackage{xurl}}{} % add URL line breaks if available
\urlstyle{same}
% fallback for those not using the hyperref driver hyperxmp:
\makeatletter
\@ifundefined{xmpquote}{\newcommand{\xmpquote}[1]{#1}}{}
\makeatother
\hypersetup{
  colorlinks=true,
  linkcolor={red},
  filecolor={red},
  citecolor={red},
  urlcolor={red},
  pdfcreator={LaTeX via pandoc}}

\title{Learning and Evaluation on the IVM Lattice}
\author{Daniel Ari Friedman\\ ORCID: 0000-0001-6232-9096\\ Email: daniel@activeinference.institute\\ DOI: 10.5281/zenodo.16887791}
\date{October 07, 2026}

\begin{document}
\maketitle

{
\hypersetup{linkcolor=black}
\setcounter{tocdepth}{3}
\tableofcontents
}
\setstretch{1.0}
\section{Learning and Evaluation on the IVM
Lattice}\label{learning-and-evaluation-on-the-ivm-lattice}

\subsection{Overview}\label{overview}

Learning on a lattice is easy to fake. A learner that merely
interpolates its observations through the graph can look excellent on
the very data it was fit to, and a dynamics model identified on early
snapshots can silently diverge on late ones. This section develops the
training/testing methodology used to evaluate the two IVM learners
honestly --- the static field learner of
\href{11_ivm_field_learning.md}{Static IVM Field Learning} and the
dynamics identifier of \href{12_ivm_dynamics.md}{IVM Lattice Dynamics}.
All utilities live in the \passthrough{\lstinline!learning\_eval.py!}
module (\passthrough{\lstinline!kfold\_site\_splits()!},
\passthrough{\lstinline!cross\_validate\_field()!},
\passthrough{\lstinline!trajectory\_train\_test()!},
\passthrough{\lstinline!learning\_curve()!},
\passthrough{\lstinline!enclosing\_radius()!}), use
\passthrough{\lstinline!numpy!} only, and are fully deterministic for
fixed seeds. Every behavioral claim below is pinned in
\passthrough{\lstinline!tests/test\_learning\_eval.py!} with fixed RNG
seeds and real numerics.

\subsection{Why held-out evaluation is harder on a
lattice}\label{why-held-out-evaluation-is-harder-on-a-lattice}

On independent data, random resampling is benign. On a lattice it is
not: field values at neighboring sites are strongly correlated, and the
learners of \passthrough{\lstinline!ivm\_field.py!} exploit exactly that
correlation --- the kernel-weighted target of the normal equations
\eqref{eq:ivm-normal-equations} predicts an unobserved site as a
weighted average of its graph neighbors. A random split that leaves
graph neighbors on both sides therefore lets the learner ``cheat'': the
held-out site is predicted from observations one hop away, and the
apparent skill measures the smoothness of the field under
\eqref{eq:ivm-objective}, not generalization to unseen territory. The
autocorrelation that makes lattice learning work is the same
autocorrelation that inflates its evaluation:

\begin{equation}
\label{eq:learn-correlation}
\operatorname{Cov}\bigl(f_i, f_j\bigr) \;\approx\; \sigma_f^{2}\,
K\bigl(d(i,j)\bigr),
\qquad d(i,j) \;=\; \text{hop distance on the IVM graph},
\end{equation}

with \(K\) the Gaussian hop kernel of \eqref{eq:ivm-kernel} and
\(\sigma_f^2\) the field variance. Because \(K\) decays with hop
distance, the leakage is local --- but the IVM ball is small, and a
random split leaks everywhere.

The methodology below therefore always holds out \emph{structure}, not
just rows: whole folds of sites for the field learner (never shared
between training and scoring), and a temporal suffix of snapshots for
dynamics identification (never shuffled). The pinned numerical result
worth internalizing: on the radius-2 ball (55 sites, 80\% observed), the
noise-free held-out error of the kernel-weighted interpolator is
dominated not by observation noise but by \emph{coverage geometry} ---
which sites happen to be observed --- so per-fold scores vary
substantially even at fixed sample size, and honest evaluation must
average over folds.

\subsection{K-fold site splits}\label{k-fold-site-splits}

\passthrough{\lstinline!kfold\_site\_splits(observed\_sites, k, seed)!}
partitions observed sites into \(k\) folds by a seeded
\passthrough{\lstinline!numpy.random.default\_rng!} permutation cut with
\passthrough{\lstinline!numpy.array\_split!} (the first \(n \bmod k\)
folds receive one extra site). Each split is a (train, test) pair; folds
are pairwise disjoint and jointly exhaustive, and site order within each
list follows the original normalized index order, so the splits are a
pure function of \((sites, k, seed)\). With \(s\) the seeded permutation
and \(F_j\) fold \(j\)'s index set, the held-out score of a fit
\(\hat{f}^{(-F_j)}_\lambda\) that never saw \(F_j\) is

\begin{equation}
\label{eq:learn-fold}
\widehat{M}(\lambda, F_j) \;=\; \frac{1}{|F_j|}
\sum_{i \in F_j}
\bigl( \hat{f}^{(-F_j)}_\lambda(i) - y_i \bigr)^{2},
\end{equation}

the quantity \passthrough{\lstinline!IVMField.score()!} returns.
Duplicate sites are rejected after normalization (two projective
representatives of one lattice point would sit on both sides of a
split), as is \(k < 2\) or \(k > n\).

\subsection{Cross-validation and honest model
selection}\label{cross-validation-and-honest-model-selection}

\passthrough{\lstinline!cross\_validate\_field(field\_values, sites, k, lam\_grid, seed)!}
runs \eqref{eq:learn-fold} for every fold and every regularization
candidate of \passthrough{\lstinline!lam\_grid!}, fitting each fold from
scratch on the training sites only --- a fresh
\passthrough{\lstinline!IVMField.lattice\_ball!} per fit, via
\passthrough{\lstinline!IVMField.learn()!} with the objective
\eqref{eq:ivm-objective}. The result is the full table
\(\widehat{M}(\lambda, F_j)\) with per-candidate mean and standard
deviation over folds, and the selected strength

\begin{equation}
\label{eq:learn-lambda-star}
\lambda^{*} \;=\; \arg\min_{\lambda \in \mathcal{G}}
\frac{1}{k} \sum_{j=1}^{k} \widehat{M}(\lambda, F_j),
\end{equation}

resolved to the earliest candidate on ties. Selection uses validation
folds only; the returned \passthrough{\lstinline!refit\_field!} is then
refit with \(\lambda^{*}\) on \emph{all} observed sites, matching the
deployment condition where every observation is available. Separating
the two is what makes the selection honest: choosing \(\lambda\) on the
same data the final model fits is how interpolation gets mistaken for
skill.

Two regimes are pinned numerically on synthetic fields. With exact
observations of a harmonic (linear in embedded XYZ) field at 80\%
coverage, cross-validation selects \(\lambda^{*} = 0\): the
kernel-weighted interpolator generalizes best, and Laplacian smoothing
only biases the fit away from exact data --- the mean held-out MSE is
strictly increasing in \(\lambda\) across a grid spanning two decades.
With noisy observations of a smooth non-harmonic field, the selection
flips to \(\lambda^{*} > 0\): interpolation carries observation noise
into the fit, and moderate regularization wins on the held-out folds.
The margin is honest about the learner's limits --- because observed
rows are pinned to their data, the noise propagates through the graph
regardless of \(\lambda\), so the achievable gain is set by how much
neighbor averaging the kernel already performs, not by the regularizer
alone.

\subsection{Temporal splits for dynamics
identification}\label{temporal-splits-for-dynamics-identification}

\passthrough{\lstinline!trajectory\_train\_test(observed, split\_frac, lattice)!}
evaluates \passthrough{\lstinline!fit\_trajectory()!}-style
identification of the coupling \(\alpha\) (\href{12_ivm_dynamics.md}{IVM
Lattice Dynamics}). The snapshot matrix (\(T+1\) rows, row 0 the initial
condition) is split \textbf{temporally}: the first
nearest-integer-rounded fraction of rows trains, the trailing rows are
held out, and no shuffling is ever applied --- a shuffled split
interpolates between adjacent snapshots of a deterministic trajectory
and hides exactly the forecast error being measured. The split point is
clamped so training keeps at least two rows (initial condition plus one
snapshot) and testing keeps at least one. Identification uses the
training prefix only, with candidate grid and refinement passed through
to \passthrough{\lstinline!fit\_trajectory()!}. Two held-out errors are
reported. The multi-step continuation error re-simulates from the
observed initial condition with the identified \(\hat{\alpha}\) and
scores only the held-out rows:

\begin{equation}
\label{eq:learn-horizon}
M_{\mathrm{test}}(\hat{\alpha}) \;=\; \frac{1}{(T - t_s)\,N}
\sum_{t=t_s+1}^{T} \big\lVert u_{\hat{\alpha}}(t) - u_{\mathrm{obs}}(t)
\big\rVert^{2} ,
\end{equation}

where \(t_s\) is the last training row. This is the stringent metric:
for any model that is not exactly right, error compounds with horizon,
so the held-out tail punishes misspecification that early rows forgive.
The one-step-ahead error is teacher-forced --- each held-out transition
is predicted a single step from the \emph{true} previous state, the
first transition starting from the last training row:

\begin{equation}
\label{eq:learn-onestep}
M_{1}(\hat{\alpha}) \;=\; \frac{1}{(T - t_s)\,N}
\sum_{t=t_s}^{T-1} \big\lVert
\mathrm{step}\bigl(u_{\mathrm{obs}}(t);\, \hat{\alpha}\bigr)
- u_{\mathrm{obs}}(t+1) \big\rVert^{2},
\end{equation}

isolating local prediction error from error accumulation. The contrast
between \eqref{eq:learn-horizon} and \eqref{eq:learn-onestep} is
diagnostic: a correct model drives both to zero (pinned exactly: a heat
trajectory with \(\alpha = 0.3\) in the candidate grid yields train,
continuation, and one-step errors of exactly \(0\)), while a
misspecified model class --- a heat model identified on
majority-dynamics data --- shows a modest training error, a held-out
continuation error more than twice as large (horizon compounding), and a
one-step error five times smaller than the continuation error (local
prediction is easier than long-horizon forecasting). Note the direction
of the misspecification: \passthrough{\lstinline!majority\_step()!}
requires integer states, so the float-continuum heat model can be fit to
integer data, but not the reverse.

\subsection{Learning curves}\label{learning-curves}

\passthrough{\lstinline!learning\_curve(values, sites, train\_fracs, seed)!}
traces the classic data-coverage curve on a lattice. One seeded
permutation of the observed sites is drawn; for each requested fraction
\(p\) the first \(m(p) = \operatorname{round}(p \cdot n)\) sites of that
permutation (clamped to \([1, n-1]\)) form the training subset and the
remainder the held-out complement:

\begin{equation}
\label{eq:learn-curve}
M(p) \;=\; \frac{1}{n - m(p)} \sum_{i \notin P_{m(p)}}
\bigl( \hat{f}_{P_{m(p)}}(i) - y_i \bigr)^{2},
\qquad P_{m} \;=\; \text{first } m \text{ entries of the permutation},
\end{equation}

reported together with the in-sample score of the same fit. The subsets
are \emph{nested} by construction --- growing \(p\) never removes an
earlier observation --- so the curve is a proper coverage curve rather
than a family of independent splits. The pinned example (radius-2 ball,
quadratic bowl observed with noise \(\sigma = 0.4\) everywhere,
\(\lambda = 0.05\)) shows held-out MSE falling by more than a factor of
two from 10\% to 85\% coverage, with non-monotone wiggles in between:
with few observations the held-out complement is large and its geometry
varies between fractions, so the curve is monotone-ish, not monotone ---
the final point sits far below the first, and that is the asserted
invariant. The curve also separates interpolation from generalization
visibly: the in-sample MSE stays below \(0.15\) at every coverage level
while the held-out MSE never drops below \(0.7\) --- the pinned
observations are fit closely at all sizes, and the difficulty is
entirely on unseen sites.

\subsection{Three-way splitting and trainable site
fits}\label{sec:learn_three_way_site_fits}

The lattice surfaces above all share one shape: a seeded split of sites
or snapshots, a fit on the training side, and scores on the held-out
side. The same shape is available for plain tabular work in the same
module (\passthrough{\lstinline!src/quadmath/learn/learning\_eval.py!}),
under the same discipline --- all split randomness confined to the
seeded permutation of the split, and fits that are pure deterministic
functions of their inputs.

\passthrough{\lstinline!three\_way\_split(n, train\_frac=0.6, val\_frac=0.2, seed=0)!}
partitions \passthrough{\lstinline!range(n)!} into disjoint
train/validation/test index lists. One seeded
\passthrough{\lstinline!numpy.random.default\_rng!} permutation of
\passthrough{\lstinline!range(n)!} is drawn and cut into three
contiguous blocks: the first \(m_{train} =
\operatorname{round}(\mathrm{train\_frac} \cdot n)\) indices train, the
next \(m_{val} = \operatorname{round}(\mathrm{val\_frac} \cdot n)\)
validate, and the remainder test. Each returned list is sorted, the
three are pairwise disjoint with union exactly
\passthrough{\lstinline!range(n)!}, and the split is a deterministic
function of \((n, \mathrm{train\_frac}, \mathrm{val\_frac}, seed)\) ---
the same deterministic-permutation contract as the k-fold site splits
above, applied to a plain index range rather than observed lattice
sites. Rounding is unguarded at the tails: for small \(n\) a trailing
block may come back empty (\(n = 3\) under the default fractions leaves
the test list empty), and a \passthrough{\lstinline!ValueError!} is
raised when either fraction is not strictly positive or their sum
reaches \(1\) --- a split with no held-out rows at all cannot be honest.

\passthrough{\lstinline!ridge\_site\_fit(features, values, lam=1e-3)!}
fits a single linear site model in closed form. With \(X_c\) the
mean-centered design matrix, \(y_c\) the mean-centered target, and
\(\bar{x}\), \(\bar{y}\) the feature-column means and target mean, the
ridge normal equations are solved exactly for the slope vector \(w\):

\begin{equation}
\label{eq:learn-ridge}
\bigl( X_c^{\!\top} X_c + \lambda I \bigr)\, w \;=\; X_c^{\!\top} y_c,
\qquad b \;=\; \bar{y} - \bar{x}^{\!\top} w,
\end{equation} with \(b\) the intercept recovered from the centered
solve; \(\lambda \geq 0\) penalizes the centered coefficient norm,
shrinking \(w\) toward zero, and \(\lambda = 0\) recovers the exact
least-squares solve. The returned \passthrough{\lstinline!RidgeSiteFit!}
carries \passthrough{\lstinline!coefficients!},
\passthrough{\lstinline!intercept!}, and the in-sample
\passthrough{\lstinline!train\_mse!}, and the fit is a deterministic
function of \passthrough{\lstinline!(features, values, lam)!} --- a
\passthrough{\lstinline!numpy.linalg.solve!}, not an iterative
optimizer, so repeated calls on the same inputs reproduce it bit for
bit. Validation is explicit: a negative \(\lambda\), a non-2-D design
matrix, a non-1-D target, disagreeing sample counts, or zero rows each
raise before any linear algebra runs.

\passthrough{\lstinline!GradientDescentTrainer(lr=0.05, max\_iters=300, tol=1e-9)!}
is the iterative counterpart, useful where the closed form is
deliberately being compared against: a full-batch gradient-descent fit
of the same linear model. \passthrough{\lstinline!fit!} standardizes the
feature columns internally (per-column zero mean and unit standard
deviation, with constant columns passing through at standard deviation
\(1\)), then runs at most \passthrough{\lstinline!max\_iters!}
full-batch updates on the mean-squared-error loss,

\begin{equation}
\label{eq:learn-gd}
\theta \;\leftarrow\; \theta \;-\; \frac{2\,\eta}{m}\,
\begin{bmatrix} X_s^{\!\top} e \\ \mathbf{1}^{\!\top} e \end{bmatrix},
\qquad e \;=\; X_s w + b - y,
\end{equation}

where \(X_s\) is the standardized design, \(m\) the row count, \(\eta\)
the learning rate, \(w\) and \(b\) the coefficient vector and bias, and
\(y\) the target. The MSE at the top of each executed iteration is
appended to \passthrough{\lstinline!loss\_history!} --- one float per
iteration, recorded before the update it describes --- and training
stops early with \passthrough{\lstinline!converged\_ = True!} as soon as
two consecutive losses differ by less than
\passthrough{\lstinline!tol!}. The fit exposes
\passthrough{\lstinline!coef\_!} expressed against the standardized
features and \passthrough{\lstinline!intercept\_!} in original target
units, together with the stored \passthrough{\lstinline!mean\_!} and
\passthrough{\lstinline!std\_!} it used, so
\passthrough{\lstinline!predict!} re-applies the recorded
standardization before the linear map and thus scores consistently
across the train, validation, and test blocks of a
\passthrough{\lstinline!three\_way\_split!}. Like
\passthrough{\lstinline!ridge\_site\_fit!} the trainer draws no
randomness of its own --- no shuffling, no stochastic gradient --- so
its entire loss history is deterministic for fixed inputs: the same rows
and hyperparameters replay the same loss curve to the last float, and
the per-iteration record doubles as the divergence/plateau diagnostic
that a scalar final loss cannot provide.

\begin{figure}
\centering
\pandocbounded{\includegraphics[keepaspectratio,alt={Convergence of the full-batch gradient-descent trainer. Per-iteration training loss (mean squared error, linear y-axis) of GradientDescentTrainer on a deterministic synthetic linear problem: a 40 \textbackslash times 3 standard-normal design matrix with true coefficients (1.5, -2.0, 0.75), intercept 0.5, and Gaussian target noise \textbackslash sigma = 0.05 (seed 0). Fitted with lr = 0.05, max\_iters = 300, tol = 1e-9; the loss falls from 6.23 at iteration 0 to the noise-floor scale \textbackslash sigma\^{}2 = 2.5\textbackslash times 10\^{}\{-3\}, flagged as converged after 126 iterations. Regenerate with quadmath/scripts/learning\_gallery.py.}]{../output/figures/learn_loss_history.png}}
\caption{\textbf{Convergence of the full-batch gradient-descent
trainer.} Per-iteration training loss (mean squared error, linear
\(y\)-axis) of \protect\passthrough{\lstinline!GradientDescentTrainer!}
on a deterministic synthetic linear problem: a \(40 \times 3\)
standard-normal design matrix with true coefficients
\((1.5, -2.0, 0.75)\), intercept \(0.5\), and Gaussian target noise
\(\sigma = 0.05\) (seed \(0\)). Fitted with
\protect\passthrough{\lstinline!lr = 0.05!},
\protect\passthrough{\lstinline!max\_iters = 300!},
\protect\passthrough{\lstinline!tol = 1e-9!}; the loss falls from
\(6.23\) at iteration \(0\) to the noise-floor scale
\(\sigma^2 = 2.5\times 10^{-3}\), flagged as converged after \(126\)
iterations. Regenerate with
\protect\passthrough{\lstinline!quadmath/scripts/learning\_gallery.py!}.}
\end{figure}

\subsection{Verification}\label{verification}

\passthrough{\lstinline!tests/test\_learning\_eval.py!} pins all of the
above with fixed seeds and no mocks: disjoint-exhaustive folds with
exact fold sizes, seed-determinism and seed-sensitivity of the splits,
the \(\lambda^{*}=0\) and \(\lambda^{*}>0\) regimes with their table
identities (mean and standard deviation recomputed from the fold table;
the refit reproduced by hand), exact-zero errors for the correctly
specified dynamics against a \(> 2\times\) continuation blow-up for the
misspecified one, manual recomputation of the one-step-ahead error,
fraction clamping at both extremes, and the learning curve's
final-below-first invariant. The module carries 100\% statement and
branch coverage, and the validation rules of every entry point (distinct
normalized sites, \(k \in
[2, n]\), fractions in the open unit interval, non-empty grids, minimum
row counts) are each exercised.

\subsection{Cross-References}\label{cross-references}

\begin{itemize}
\tightlist
\item
  Coordinate foundations and the twelve-around-one shell:
  \href{03_quadray_methods.md}{Quadray Methods}
\item
  The learner being evaluated: \href{11_ivm_field_learning.md}{Static
  IVM Field Learning}
\item
  The dynamics being identified: \href{12_ivm_dynamics.md}{IVM Lattice
  Dynamics}
\item
  Enumeration and search infrastructure:
  \href{13_lattice_tooling.md}{Lattice Tooling}
\end{itemize}

\end{document}
